% theme: applications_of_prices_and_quantities
\MajorQuestionLesson{代金と個数の利用問題}
\SetRowSpace{6mm}

\TypeQuestionDouble{品物の個数を $x,\ y$ 個として連立方程式をつくり、それぞれの個数を求めなさい。}{kosu}

\QQ{1本80円の鉛筆と1個120円の消しゴムを合わせて15個買うと、代金の合計は1560円であった。鉛筆と消しゴムをそれぞれ何個買ったか。}{}
\QQ{1冊150円のノートと1本90円のボールペンを合わせて18個買うと、代金の合計は2340円であった。ノートとボールペンをそれぞれ何個買ったか。}{}
\QQ{1個70円のパンと1本110円の飲み物を合わせて24個買うと、代金の合計は2200円であった。パンと飲み物をそれぞれ何個買ったか。}{}
\QQ{1個180円のケーキと1個250円のプリンを合わせて14個買うと、代金の合計は2940円であった。ケーキとプリンをそれぞれ何個買ったか。}{}

\TypeQuestionDouble{求める2つの数量をそれぞれ $x,\ y$ として連立方程式をつくり、問いに答えなさい。}{ryoukin}

\QQ{ある施設の入場料は、大人1人700円、子ども1人400円である。大人と子どもが合わせて23人入場し、入場料の合計は12200円であった。大人と子どもはそれぞれ何人か。}{}
\QQ{100円硬貨と50円硬貨が合わせて32枚あり、金額の合計は2550円である。100円硬貨と50円硬貨はそれぞれ何枚あるか。}{}
\QQ{1個600円の弁当を、閉店前には1個400円で販売した。合わせて28個売れ、売上金額は14800円であった。600円で売れた弁当と400円で売れた弁当は、それぞれ何個か。}{}
\QQ{ある映画館の入場料は、大人1人1200円、学生1人800円である。大人と学生が合わせて37人入場し、入場料の合計は35600円であった。大人と学生はそれぞれ何人か。}{}

\LessonPageBreak

\TypeQuestionDouble{品物1個の値段を $x,\ y$ 円として連立方程式をつくり、それぞれの値段を求めなさい。}{nedan}

\QQ{同じ種類のノート3冊とボールペン2本を買うと570円、ノート2冊とボールペン5本を買うと710円であった。ノート1冊とボールペン1本の値段をそれぞれ求めなさい。}{}
\QQ{同じ種類のりんご4個とみかん3個を買うと720円、りんご2個とみかん5個を買うと640円であった。りんご1個とみかん1個の値段をそれぞれ求めなさい。}{}
\QQ{同じ種類のサンドイッチ2個と飲み物3本を買うと1120円、サンドイッチ5個と飲み物2本を買うと2030円であった。サンドイッチ1個と飲み物1本の値段をそれぞれ求めなさい。}{}
\QQ{ある施設で、大人4人と学生7人の入場料は合わせて7800円、大人9人と学生3人の入場料は合わせて9900円であった。大人1人と学生1人の入場料をそれぞれ求めなさい。}{}

\TypeQuestionDouble{求める2つの値段をそれぞれ $x,\ y$円 として連立方程式をつくり、問いに答えなさい。}{otsuri}

\QQ{同じ種類の鉛筆5本とノート2冊を買い、1000円を出すと140円のおつりがあった。また、鉛筆3本とノート4冊を買い、1200円を出すと180円のおつりがあった。鉛筆1本とノート1冊の値段をそれぞれ求めなさい。}{}
\QQ{同じ種類のケーキ3個とプリン4個を買い、2000円を出すと560円のおつりがあった。また、ケーキ5個とプリン2個を買おうとすると、1500円では60円不足した。ケーキ1個とプリン1個の値段をそれぞれ求めなさい。}{}
\QQ{同じ種類の赤いペン4本と青いペン3本を買い、50円の包装代を加えると1010円であった。また、赤いペン2本と青いペン5本を買い、50円の包装代を加えると1090円であった。赤いペン1本と青いペン1本の値段をそれぞれ求めなさい。}{}
\QQ{同じ種類の商品Aを4個と商品Bを3個注文すると、400円の送料を含めて1730円であった。また、商品Aを2個と商品Bを6個注文すると、400円の送料を含めて1740円であった。商品A、商品Bの1個の値段をそれぞれ求めなさい。}{}
